\chapter[Zasady maksimum i teoria potencjału]{Zasady maksimum, lemat Hopfa, nierówności Harnacka i podstawy teorii potencjału}
\label{chap:pde14}

W rozdziale~\ref{chap:pde13} uzyskaliśmy rozwiązania słabe
zagadnień eliptycznych metodami wariacyjnymi. Samo istnienie w
$H^1$ nie mówi jeszcze, jaki znak ma rozwiązanie, czy osiąga ono
maksimum wewnątrz obszaru ani jak wartości w różnych punktach
zależą od siebie. Odpowiedzi na te pytania dostarczają zasady
maksimum i nierówności Harnacka. Zaczniemy od współczynników
mierzalnych, dla których wystarczy argument energetyczny, a potem
przejdziemy do silniejszych twierdzeń dla funkcji harmonicznych.

\section{Słaba zasada maksimum dla rozwiązań słabych}
\label{sec:pde14-weak}
Niech $\Omega\subset\R^d$ będzie obszarem ograniczonym Lipschitzowskim. Załóżmy, że $A\in L^\infty(\Omega;\R^{d\times d})$
spełnia $\xi^{\mathsf T}A(x)\xi\ge\lambda|\xi|^2$
dla pewnego $\lambda>0$, prawie każdego $x\in\Omega$ i każdego
$\xi\in\R^d$. Nie zakładamy symetrii $A$.
Rozważamy operator $Lu=-\diver(A\nabla u)+cu$, gdzie
$c\in L^\infty(\Omega)$ i $c\ge0$ prawie wszędzie.
Wszystkie funkcje i współczynniki są rzeczywiste.
Stosujemy konwencje rozdziału~\ref{chap:pde13}: nierówność
eliptyczności zachodzi na wspólnym zbiorze pełnej miary,
a formę $a$ definiujemy na $H^1(\Omega)\times H^1(\Omega)$.

\begin{definicja}[Słabe rozwiązanie nierówności]
\label{def:pde14-subsolution}
Funkcja $u\in H^1(\Omega)$ jest \emph{słabym podrozwiązaniem}
$Lu\le0$ w $\Omega$, jeżeli
\[
 \int_\Omega\big(A\nabla u\cdot\nabla\varphi
 +cu\varphi\big)\,\mathrm d\lambda_d(x)\le0
 \quad\text{dla każdego }\varphi\in H_0^1(\Omega),\ \varphi\ge0.
\]
Analogicznie definiujemy nadrozwiązanie, odwracając nierówność.
Całka istnieje i jest ciągła względem $\varphi$ w normie
$H^1$, na mocy oszacowania formy~\eqref{eq:pde13-form}.
Jeżeli $u$ jest jednocześnie podrozwiązaniem i nadrozwiązaniem,
to $a(u,\varphi)=0$ dla nieujemnych $\varphi\in H_0^1$.
Dla dowolnego rzeczywistego $\varphi\in H_0^1$ piszemy
$\varphi=\varphi^+-\varphi^-$; obie części należą do $H_0^1$
na mocy wniosku~\ref{cor:pde9-truncation},
lematu~\ref{lem:pde13-trace-parts} i twierdzenia
~\ref{thm:pde10-trace-zero}. Z liniowości otrzymujemy
równość dla wszystkich testów, czyli rozwiązanie słabe
$Lu=0$ (bez narzucania w tej definicji danych brzegowych).
Odwrotna implikacja wynika przez ograniczenie do testów nieujemnych.
W definicji nie zakładamy istnienia klasycznych pochodnych drugiego rzędu.
\end{definicja}

\begin{lemat}[Obcięcie poziomowe]
\label{lem:pde14-trunc}
Jeżeli $u\in H^1(\Omega)$ i $k\in\R$, to
$w=(u-k)^+:=\max\{u-k,0\}\in H^1(\Omega)$ oraz
$\nabla w=\chi_{\{u>k\}}\nabla u$ prawie wszędzie.
Jeżeli ponadto $\operatorname{Tr}u\le k$ prawie wszędzie na
$\partial\Omega$, to $w\in H_0^1(\Omega)$.
\end{lemat}
\begin{proof}
Ponieważ $\Omega$ jest ograniczona, stała $k$ należy do
$H^1(\Omega)$, ma zerowy słaby gradient i ślad równy $k$.
Stosujemy wniosek~\ref{cor:pde9-truncation} z $p=2$
do funkcji $u-k$. Otrzymujemy $w\in H^1$ i
$\nabla w=\chi_{\{u-k>0\}}\nabla(u-k)=\chi_{\{u>k\}}\nabla u$
prawie wszędzie.
Z lematu~\ref{lem:pde13-trace-parts} oraz liniowości śladu
wynika $\operatorname{Tr}w=(\operatorname{Tr}u-k)^+$.
Przy założeniu na ślad prawa strona znika.
Twierdzenie~\ref{thm:pde10-trace-zero} daje $w\in H_0^1$.
Nie powtarzamy dowodów obcięcia i zgodności śladu z częściami
dodatnimi z wcześniejszych rozdziałów.
\end{proof}

Poniższy wynik uogólnia elementarną zasadę maksimum
z twierdzenia~\ref{thm:laplace-max}: tutaj współczynniki
są jedynie mierzalne, a równanie i warunek brzegowy
interpretujemy w sensie słabym. Hipotezy obu twierdzeń
są różne, więc nie utożsamiamy ich dowodów.

\begin{twierdzenie}[Słaba zasada maksimum]
\label{thm:pde14-weak-max}
Niech $u\in H^1(\Omega)$ będzie podrozwiązaniem $Lu\le0$,
przy powyższych założeniach o $A$ i $c$. Niech $M\ge0$ będzie
liczbą taką, że $\operatorname{Tr}u\le M$ prawie wszędzie na
$\partial\Omega$. Wówczas $u\le M$ prawie wszędzie w $\Omega$.
Jeżeli $c=0$, ten sam wniosek jest prawdziwy dla dowolnego
$M\in\R$ spełniającego warunek na ślad.
\end{twierdzenie}
\begin{proof}
Kładziemy $w=(u-M)^+$. Z lematu~\ref{lem:pde14-trunc}
$w\in H_0^1(\Omega)$ i możemy użyć go jako nieujemnej
funkcji testowej. Na zbiorze $\{u>M\}$ mamy
$\nabla w=\nabla u$ i $u=M+w$, natomiast poza nim
$w=\nabla w=0$ prawie wszędzie. Zatem
\[
 0\ge \int_\Omega(A\nabla u\cdot\nabla w+cuw)\,
 \mathrm d\lambda_d(x)
 \ge \lambda\|\nabla w\|_2^2
      +\int_{\{u>M\}}c(M+w)w\,\mathrm d\lambda_d(x).
\]
Ponieważ $c\ge0$, $M\ge0$ i $w\ge0$, ostatnia całka jest
nieujemna. Ponieważ $\lambda>0$, mamy $\|\nabla w\|_2=0$,
a nierówność Poincarégo dla $H_0^1(\Omega)$
(wniosek~\ref{cor:pde10-Poincare}) daje $\|w\|_2=0$,
czyli $u\le M$ prawie wszędzie. Jeśli $c=0$, składnik
zerowego rzędu znika, więc znak $M$ nie odgrywa roli.
\end{proof}

\begin{wniosek}[Porównanie]
\label{cor:pde14-compare}
Niech $u,v\in H^1(\Omega)$, $c\ge0$, a
$L=-\diver(A\nabla\cdot)+c$. Jeżeli $Lu\le Lv$ słabo w
$\Omega$ i $\operatorname{Tr}u\le\operatorname{Tr}v$ prawie
wszędzie na $\partial\Omega$, to $u\le v$ prawie wszędzie
w $\Omega$.
\end{wniosek}
\begin{proof}
Różnica $w=u-v$ jest podrozwiązaniem $Lw\le0$ i ma
niedodatni ślad. Stosujemy twierdzenie~\ref{thm:pde14-weak-max}
z $M=0$. Nie wykorzystujemy symetrii macierzy $A$.
\end{proof}

\begin{uwaga}[Granice wniosku]
Przy członie $b\cdot\nabla u$ poprzedni test zawiera
$\int_\Omega(b\cdot\nabla w)w\,\mathrm d\lambda_d(x)$,
którego znak nie jest określony. Dodatkowy argument
energetyczny podano już w twierdzeniu
~\ref{thm:pde13-drift-comparison}: dla $b\in L^\infty$
o dystrybucyjnej dywergencji w $L^\infty$ wystarcza
$c-\tfrac12\diver b\ge0$ do porównania.
Jeśli dodatkowo $c\ge0$ i $M\ge0$, stała $M$ jest
nadrozwiązaniem, gdyż $LM=cM\ge0$; porównanie daje
wtedy również $u\le M$ dla podrozwiązania $Lu\le0$.
Bez dalszych założeń nie wolno pomijać warunku znakowego
na $c$, co pokazuje poniższy przykład.
\end{uwaga}

\begin{kontrprzyklad}[Ujemny człon zerowego rzędu]
\label{ex:pde14-negative-c}
Na $\Omega=(0,1)$ funkcja $u(x)=\sin(\pi x)$ jest dodatnia
we wnętrzu i ma zerowy ślad, choć spełnia
$-u''-\pi^2u=0$. Gdyby dla $c=-\pi^2$ zachodziła zasada
maksimum z twierdzenia~\ref{thm:pde14-weak-max},
wynikałoby $u\le0$, co jest fałszem. Jest to jednocześnie
przykład rezonansu z rozdziału~\ref{chap:pde13}.
\end{kontrprzyklad}

\section{Funkcje harmoniczne i mocna zasada maksimum}
\label{sec:pde14-strong}
Od tej sekcji przyjmujemy $d\ge2$, chyba że zaznaczono inaczej.
Funkcją harmoniczną w otwartym zbiorze $U\subset\R^d$
nazywamy rzeczywistą funkcję $u\in C^2(U)$ spełniającą
$\Delta u=0$. Z lematu Weyla~\ref{thm:pde7-weyl} wiemy,
że każda rzeczywista dystrybucja $T$ spełniająca $\Delta T=0$
jest dystrybucją generowaną przez jedyną gładką funkcję
harmoniczną. Jest ona rzeczywista, ponieważ część urojona
generuje dystrybucję zerową, a stwierdzenie
~\ref{prop:pde7-injective} zapewnia jej znikanie prawie
wszędzie; z ciągłości wynika znikanie wszędzie.
Twierdzenia punktowe stosujemy do tego reprezentanta.
W wymiarze jeden funkcje harmoniczne na przedziałach są
afiniczne ($u''=0$), więc podstawowe zasady maksimum
wynikają bezpośrednio z tej postaci.

\begin{twierdzenie}[Mocna zasada maksimum dla laplasjanu]
\label{thm:pde14-strong}
Niech $U\subset\R^d$ będzie obszarem, a $u\in C^2(U)$ niech spełnia $\Delta u=0$ w $U$.
Jeżeli $u$ osiąga globalne maksimum lub globalne minimum
w pewnym punkcie $x_0\in U$, to $u$ jest stała na $U$.
Silniejszą wersję dla ekstremum lokalnego udowodniono
w twierdzeniu~\ref{thm:pde3-strong-max}.
\end{twierdzenie}
\begin{proof}
Załóżmy, że $u(x_0)=M=\sup_U u$. Dla dowolnego
$r>0$ takiego, że $\overline{B_r(x_0)}\subset U$, własność
wartości średniej (twierdzenie~\ref{thm:pde3-mean}) daje
\[
 M=u(x_0)=\frac{1}{\sigma(\partial B_r(x_0))}
 \int_{\partial B_r(x_0)}u(y)\,\mathrm d\sigma(y),
\]
gdzie $\sigma$ jest miarą powierzchniową sfery.
Funkcja $M-u$ jest nieujemna i ciągła na sferze.
Jeżeli byłaby dodatnia w choć jednym jej punkcie, to
byłaby dodatnia na pewnym otwartym fragmencie sfery,
a średnia $u$ byłaby ściśle mniejsza od $M$.
Tak więc $u\equiv M$ na każdej takiej sferze, a zatem
również w całej małej kuli.
Zbiór $E=\{x\in U:u(x)=M\}$ jest niepusty i domknięty
w $U$ z ciągłości. Ten sam argument zastosowany do każdego
$x\in E$ pokazuje, że $E$ jest też otwarty w $U$.
Spójność $U$ wymusza $E=U$. Dla minimum stosujemy wynik do $-u$.
\end{proof}

\begin{wniosek}[Dodatniość i jednoznaczność]
\label{cor:pde14-positive}
Jeżeli $U$ jest spójna, a $u\ge0$ jest harmoniczna w $U$,
to $u\equiv0$ albo $u(x)>0$ dla wszystkich $x\in U$.
Jeżeli $U$ jest ponadto ograniczona oraz
$u,v\in C^2(U)\cap C(\overline U)$ są harmoniczne i mają
te same wartości na $\partial U$, to $u=v$ na $U$.
\end{wniosek}
\begin{proof}
Jeżeli $u(x_0)=0$ we wnętrzu, funkcja osiąga tam minimum,
więc jest stała, a jej wartość w $x_0$ wynosi zero,
zatem $u\equiv0$.
Dla jednoznaczności $w=u-v$ jest harmoniczna i znika na
brzegu. Ciągłość na zwartym $\overline U$ daje istnienie
maksimum i minimum; gdyby którakolwiek ekstremalna
wartość była różna od zera, ekstremum osiągnięte byłoby
we wnętrzu, co w myśl twierdzenia~\ref{thm:pde14-strong}
wymusiłoby stałość $w$ na $U$. Ponieważ niepusty
ograniczony obszar ma niepusty brzeg, ciągłość $w$
na $\overline U$ i zerowe wartości brzegowe wymuszają
wtedy stałą równą zero, sprzecznie z niezerowym ekstremum.
Stąd oba ekstrema są zerowe, czyli $u=v$.
\end{proof}

\section{Lemat Hopfa: znak pochodnej normalnej}
\label{sec:pde14-hopf}
Słaba zasada maksimum mówi, gdzie może znajdować się
minimum. Lemat Hopfa wyjaśnia, jak funkcja zachowuje się
przy zbliżaniu do punktu minimum na brzegu.

\begin{definicja}[Warunek kuli wewnętrznej w punkcie]
Niech $U\subset\R^d$ będzie otwarta, a $x_0\in\partial U$.
Mówimy, że w $x_0$ zachodzi \emph{warunek kuli wewnętrznej},
jeżeli istnieją $z\in U$ i $R>0$ takie, że
$B_R(z)\subset U$ oraz $x_0\in\partial B_R(z)$.
Wektor $\nu=(x_0-z)/R$ jest zewnętrzną normalną do tej kuli.
\end{definicja}

\begin{twierdzenie}[Lemat Hopfa dla funkcji harmonicznych]
\label{thm:pde14-hopf}
Niech $U\subset\R^d$ będzie obszarem ograniczonym,
$x_0\in\partial U$ i niech w $x_0$ zachodzi warunek kuli
wewnętrznej $B_R(z)\subset U$. Załóżmy, że
$u\in C^2(U)\cap C(\overline U)$ jest harmoniczna, nie jest
stała, a $u(x_0)=\min_{\overline U}u$.
Wówczas dla kierunku $-\nu=(z-x_0)/R$ zachodzi
\[
 \liminf_{t\downarrow0}\frac{u(x_0-t\nu)-u(x_0)}{t}>0.
\]
Jeżeli dodatkowo istnieje skończona granica
\[
 \partial_\nu u(x_0):=
 \lim_{t\downarrow0}\frac{u(x_0)-u(x_0-t\nu)}{t},
\]
to $\partial_\nu u(x_0)<0$. Jest to pochodna w kierunku
normalnej do kuli, obliczana od strony wnętrza; gdy $u$
ma przedłużenie klasy $C^1$ do otoczenia $x_0$,
jest równa $\nabla u(x_0)\cdot\nu$.
Nie zakładamy istnienia normalnej do samego $\partial U$.
\end{twierdzenie}
\begin{proof}
Odejmując stałą, zakładamy $u(x_0)=0$ i $u\ge0$ na
$\overline U$. Z mocnej zasady minimum
$u>0$ wewnątrz $U$. W pierścieniu
$\mathcal A=\{x:R/2<|x-z|<R\}$ rozważmy funkcję
\[
 w(x)=e^{-\alpha|x-z|^2}-e^{-\alpha R^2},\qquad \alpha>0.
\]
Jest ona dodatnia wewnątrz pierścienia, a na jego zewnętrznej
sferze wynosi zero. Różniczkowanie daje dla $j=1,\ldots,d$
\[
 \partial_j w=-2\alpha(x_j-z_j)e^{-\alpha|x-z|^2},
 \qquad
 \partial_{jj}w=(4\alpha^2(x_j-z_j)^2-2\alpha)
 e^{-\alpha|x-z|^2}.
\]
Po zsumowaniu $\Delta w=(4\alpha^2|x-z|^2-2d\alpha)
e^{-\alpha|x-z|^2}$. Jeśli $\alpha>2d/R^2$, to dla
$|x-z|\ge R/2$ mamy
$4\alpha^2|x-z|^2-2d\alpha\ge\alpha(\alpha R^2-2d)>0$.
Stąd $\Delta w>0$ na $\overline{\mathcal A}$. Z ciągłości i dodatniości $u$
na zwartej wewnętrznej sferze
$m:=\min_{|x-z|=R/2}u(x)>0$. Wybieramy
$0<\varepsilon\le m/\max_{|x-z|=R/2}w(x)$.
Wtedy $u-\varepsilon w\ge0$ na obu składowych brzegu
pierścienia: na wewnętrznej z wyboru $\varepsilon$, na
zewnętrznej z $u\ge0$ oraz $w=0$. Ponieważ
$\Delta(u-\varepsilon w)=-\varepsilon\Delta w<0$,
funkcja $u-\varepsilon w$ nie może osiągnąć ujemnego
minimum we wnętrzu $\mathcal A$: w takim punkcie hesjan
byłby dodatnio półokreślony i laplasjan nieujemny.
Ciągłość na zwartym $\overline{\mathcal A}$ oraz nieujemne
wartości brzegowe dają więc $u\ge\varepsilon w$.
Dla $0<t<R/2$ punkt $x=x_0-t\nu$ leży w $\mathcal A$,
a $|x-z|=R-t$. Otrzymujemy
\[
 \frac{u(x_0-t\nu)}t\ge
 \varepsilon\frac{e^{-\alpha(R-t)^2}-e^{-\alpha R^2}}t
 \longrightarrow 2\varepsilon\alpha R e^{-\alpha R^2}>0.
\]
W konsekwencji granica dolna lewej strony jest nie mniejsza
niż $2\varepsilon\alpha R e^{-\alpha R^2}>0$.
Jeśli istnieje pochodna zdefiniowana w tezie, to lewa strona
zbiega do $-\partial_\nu u(x_0)$ (po przywróceniu odjętej
wartości $u(x_0)$), co daje tezę.
\end{proof}

\begin{przyklad}[Konieczność geometrii brzegu]
\label{ex:pde14-cusp}
W klinie płaskim $U=\{(r\cos\theta,r\sin\theta):
0<r<1,\ 0<\theta<\omega\}$, gdzie $0<\omega<\pi$,
funkcja $u(r,\theta)=r^{\pi/\omega}\sin(\pi\theta/\omega)$
jest dodatnia we wnętrzu, ponieważ $0<\pi\theta/\omega<\pi$.
Niech $p=\pi/\omega>1$. Ze wzoru laplasjanu we współrzędnych
biegunowych otrzymujemy
\[
 \Delta u=\partial_{rr}u+r^{-1}\partial_r u
 +r^{-2}\partial_{\theta\theta}u
 =[p(p-1)+p-p^2]r^{p-2}\sin(p\theta)=0.
\]
Funkcja przedłuża się ciągle na domknięcie (w zerze
$|u|\le r^p\to0$), jest nieujemna na brzegu i znika
na obu ramionach. Przy wierzchołku ma minimum, lecz
$0\le u(r,\theta)/r\le r^{p-1}\to0$, nawet jednostajnie
w kierunku $\theta$.
W zerze nie zachodzi warunek kuli wewnętrznej.
Gdyby $B_R(z)\subset U$ i $|z|=R$, to środek $z$
o kącie $\theta_z\in(0,\omega)$ musiałby mieć odległość
od każdej z prostych ograniczających klin co najmniej $R$,
ponieważ klin jest zawarty w przecięciu odpowiednich
półpłaszczyzn. Odległości te wynoszą
$R\sin\theta_z$ i $R\sin(\omega-\theta_z)$.
Co najmniej jeden z kątów $\theta_z,\omega-\theta_z$
jest mniejszy niż $\pi/2$, więc jedna odległość jest
ściśle mniejsza niż $R$, sprzeczność.
Nie wolno zatem stosować lematu Hopfa z dodatnią liniową
dolną granicą przy wierzchołku.
\end{przyklad}

\section{Jądro Poissona i ilościowe porównywanie wartości}
\label{sec:pde14-poisson}
Oznaczmy przez $\omega_{d-1}$ pole powierzchni
jednostkowej sfery $S^{d-1}$ dla $d\ge2$ (w rozdziale
~\ref{chap:pde-laplace} tę samą wielkość oznaczono przez
$\omega_d$). Przez $\sigma$ rozumiemy miarę powierzchniową.
W kuli
$B_R=B_R(0)\subset\R^d$ jądro Poissona ma postać
\[
 P_R(x,y)=\frac{R^2-|x|^2}{\omega_{d-1}R|x-y|^d},
 \qquad x\in B_R,\quad y\in\partial B_R.
\]
Jest dodatnie i spełnia
$\int_{\partial B_R}P_R(x,y)\,\mathrm d\sigma(y)=1$.
Dla $g\in C(\partial B_R)$ funkcja
$u(x)=\int_{\partial B_R}P_R(x,y)g(y)\,\mathrm d\sigma(y)$
jest harmoniczna i przyjmuje dane brzegowe $g$ w sensie
ciągłego przedłużenia na $\overline{B_R}$; są to wyniki
twierdzenia~\ref{thm:pde3-poisson}. Jeżeli $u$ jest już
harmoniczna w $B_R$ i ciągła na $\overline{B_R}$,
stosujemy ten wynik do $g=u|_{\partial B_R}$.
Różnica między $u$ a całką Poissona jest harmoniczna
i ma zerowe wartości brzegowe, więc znika na mocy
wniosku~\ref{cor:pde14-positive}. Otrzymujemy formułę
Poissona dla dowolnej takiej funkcji $u$.

\begin{twierdzenie}[Nierówność Harnacka w kuli]
\label{thm:pde14-harnack}
Niech $d\ge2$, $0<r<R$, a $u\ge0$ będzie funkcją
harmoniczną w otoczeniu $\overline{B_R(0)}$.
Wtedy dla każdego $x\in\overline{B_r(0)}$ zachodzi
\begin{equation}\label{eq:pde14-harnack}
 \frac{R^2-r^2}{R^2}\left(\frac{R}{R+r}\right)^d u(0)
 \le u(x)\le
 \frac{R^{d-2}(R+r)}{(R-r)^{d-1}}u(0).
\end{equation}
W szczególności $\sup_{B_r}u\le C_{d,r/R}\inf_{B_r}u$
przy stałej zależnej wyłącznie od $d$ i $r/R$.
\end{twierdzenie}
\begin{proof}
Stosujemy stwierdzenie~\ref{prop:pde3-harnack} z $n=d$;
jego założenia wynikają z harmoniczności w otoczeniu
domknięcia kuli i nieujemności $u$. Przypomnijmy rachunek
stałych. Dla $s=|x|\le r$ i $|y|=R$ mamy
$R-s\le|x-y|\le R+s$ oraz
\[
 \frac{P_R(x,y)}{P_R(0,y)}
 =\frac{R^{d-2}(R-s)(R+s)}{|x-y|^d}.
\]
Dolna granica ilorazu wynosi co najmniej
$R^{d-2}(R-s)/(R+s)^{d-1}$, a górna nie przekracza
$R^{d-2}(R+s)/(R-s)^{d-1}$. Pierwsza funkcja $s$
jest malejąca, druga rosnąca na $[0,R)$, ponieważ
odpowiednio licznik maleje, a mianownik rośnie, lub
odwrotnie. Zastępujemy zatem $s$ przez $r$.
Mnożymy nierówności przez $P_R(0,y)u(y)\ge0$
i całkujemy względem $\mathrm d\sigma(y)$ na $\partial B_R$.
Formuła Poissona daje~\eqref{eq:pde14-harnack}; dolna
stała jest równa $R^{d-2}(R-r)/(R+r)^{d-1}$.
Oznaczmy dolną i górną stałą przez $a_r>0$ i $b_r>0$.
Wtedy $\inf_{B_r}u\ge a_ru(0)$ oraz
$\sup_{B_r}u\le b_ru(0)$, więc bez dzielenia przez
wartość funkcji otrzymujemy
\[
 \sup_{B_r}u\le\frac{b_r}{a_r}\inf_{B_r}u
 =\left(\frac{R+r}{R-r}\right)^d\inf_{B_r}u.
\]
Możemy wziąć $C_{d,r/R}=((1+r/R)/(1-r/R))^d$.
W szczególności dla $r=R/2$ otrzymujemy stałą $3^d$.
\end{proof}

\begin{uwaga}[Usunięcie założenia o przedłużeniu]
Jeżeli $u\ge0$ jest harmoniczna jedynie w $B_R$, bez
ciągłości na $\partial B_R$, stosujemy twierdzenie w
kulach $B_\rho$ o $r<\rho<R$, gdyż ich domknięcia
są zawarte w $B_R$. Dla ustalonego $r$ i $x\in\overline{B_r}$
wartości $u(0),u(x)$ nie zależą od $\rho$. Przechodząc
do granicy $\rho\uparrow R$ po obu stronach oszacowania,
korzystamy tylko z ciągłości jawnych stałych dla $\rho>r$.
Otrzymujemy te same nierówności bez żadnego założenia
o istnieniu wartości brzegowych na $\partial B_R$.
\end{uwaga}

\begin{wniosek}[Nierówność Harnacka na zwartych podzbiorach]
\label{cor:pde14-harnack-chain}
Niech $U\subset\R^d$ będzie obszarem,
$K\subset U$ będzie zwarty i niepusty, a $u\ge0$ harmoniczna
w $U$. Istnieje stała $C=C(U,K,d)>0$, niezależna od $u$,
taka że $\sup_Ku\le C\inf_Ku$.
\end{wniosek}
\begin{proof}
Wybieramy punkt $x_*\in K$. Zbiór punktów $U$, które
można połączyć z $x_*$ łamaną zawartą w $U$, jest niepusty,
otwarty i względnie domknięty: dla dowolnego punktu $z\in U$
istnieje kula o środku $z$ zawarta w $U$, a dowolne dwa
jej punkty można połączyć odcinkiem. Jeśli jeden punkt
kuli jest osiągalny łamaną, osiągalne są wszystkie.
Ze spójności $U$ każdy $x\in K$ można więc połączyć
z $x_*$ taką łamaną $\gamma:[0,1]\to U$.

Obraz $\gamma([0,1])$ jest zwarty. Istnieje $\rho>0$
takie, że $\overline{B_{2\rho}(z)}\subset U$ dla każdego
$z$ z obrazu łamanej: jeśli $\R^d\setminus U$ jest
niepusty, wybieramy $2\rho$ mniejsze od dodatniej odległości
tego zwartego zbioru od domkniętego dopełnienia $U$;
dla $U=\R^d$ wystarczy dowolne $\rho>0$.
Z jednostajnej ciągłości $\gamma$ wybieramy podział
$0=t_0<\cdots<t_N=1$ taki, że
$|z_{j+1}-z_j|<\rho$, gdzie $z_j=\gamma(t_j)$.
Mamy $z_0=x_*$, $z_N=x$.
Nierówność Harnacka w każdej kuli $B_{2\rho}(z_j)$
z promieniem wewnętrznym $\rho$ porównuje $u(z_j)$
i $u(z_{j+1})$ w obu kierunkach ze stałą $3^d$.
Iterując, otrzymujemy obie nierówności
$u(x)\le 3^{dN}u(x_*)$ i $u(x_*)\le 3^{dN}u(x)$.
Ostatnia kula $B_{2\rho}(x)$ pozwala porównać jeszcze
$u(x)$ z $u(y)$ dla każdego $y\in B_\rho(x)$.
Zatem dla wszystkich takich $y$ zachodzi
\[
 C_x^{-1}u(x_*)\le u(y)\le C_xu(x_*),
 \qquad C_x=3^{d(N+1)}.
\]
Kule $B_\rho(x)$, z promieniem zależnym od wybranej łamanej,
pokrywają $K$. Wybieramy skończone podpokrycie i oznaczamy
przez $C_0$ maksimum odpowiadających stałych $C_x$.
Wszystkie konstrukcje zależą od $U,K,d$, a nie od $u$.
Dla dowolnych $y,z\in K$ mamy
$u(y)\le C_0u(x_*)\le C_0^2u(z)$.
Biorąc supremum po $y$ i infimum po $z$ otrzymujemy tezę
z $C=C_0^2$, również gdy $u(x_*)=0$.
\end{proof}

\begin{wniosek}[Twierdzenie Liouville'a z ograniczeniem jednostronnym]
\label{cor:pde14-liouville}
Jeżeli $u$ jest harmoniczna na $\R^d$, $d\ge2$, i jest
ograniczona z dołu albo z góry, to $u$ jest stała.
\end{wniosek}
\begin{proof}
Jeżeli $u\ge m$ dla pewnego $m\in\R$, zastępujemy
$u$ przez $u-m$, która jest nadal harmoniczna i nieujemna.
Wystarczy więc udowodnić wynik dla $u\ge0$.
Dla ustalonego $x\in\R^d$ wybieramy $R>|x|$ i stosujemy
\eqref{eq:pde14-harnack} z $r=|x|$ (dla $x=0$ teza jest
oczywista). Gdy $R\to\infty$, oba współczynniki przy $u(0)$
dążą do $1$. Z nierówności i przejścia do granicy otrzymujemy
$u(x)=u(0)$. Jeśli pierwotna funkcja jest ograniczona
z góry przez $M$, stosujemy dowód do $M-u\ge0$.
W obu przypadkach pierwotna funkcja jest stała.
Jest to wzmocnienie wniosku~\ref{cor:pde3-liouville},
w którym zakładano ograniczenie obustronne.
\end{proof}

\section{Oszacowania pochodnych i lokalne wygładzanie}
\label{sec:pde14-deriv}
\begin{twierdzenie}[Oszacowanie gradientu funkcji harmonicznej]
\label{thm:pde14-gradient}
Niech $d\ge2$, $R>0$, $x_0\in\R^d$ i niech $u$ będzie
harmoniczna w otoczeniu $\overline{B_R(x_0)}$.
Wtedy
\[
 |\nabla u(x_0)|\le \frac dR
 \sup_{B_R(x_0)}|u|.
\]
W szczególności, jeżeli $0<r<R$ i $u$ jest harmoniczna
w $B_R(x_0)$, to
$\sup_{B_r(x_0)}|\nabla u|\le
\frac{2d}{R-r}\sup_{B_R(x_0)}|u|$.
\end{twierdzenie}
\begin{proof}
Przez translację możemy przyjąć $x_0=0$.
Uzupełniamy składowe oszacowanie ze stwierdzenia
~\ref{prop:pde3-gradient} do oszacowania normy euklidesowej
całego gradientu. Różniczkujemy formułę Poissona:
na $\overline{B_{R/2}}\times\partial B_R$ mamy $|x-y|\ge R/2$,
więc wszystkie pierwsze pochodne jądra są ciągłe
i ograniczone. Ilorazy różnicowe względem każdej współrzędnej
w punkcie $0$ są jednostajnie ograniczone przez supremum
odpowiedniej pochodnej (twierdzenie o wartości średniej).
Po pomnożeniu przez ograniczoną funkcję $u(y)$ otrzymujemy
majorantę całkowalną na sferze. Twierdzenie o zbieżności
zmajoryzowanej uzasadnia różniczkowanie pod całką.
Rachunek daje
\[
 \nabla_x P_R(x,y)=\frac{1}{\omega_{d-1}R}
 \left[-2x|x-y|^{-d}
 -d(R^2-|x|^2)(x-y)|x-y|^{-d-2}\right].
\]
Dla $|y|=R$ mamy zatem
$P_R(0,y)=1/(\omega_{d-1}R^{d-1})$ oraz
$\nabla_x P_R(0,y)=d y/(\omega_{d-1}R^{d+1})$. Stąd
\[
 \nabla u(0)=\frac{d}{\omega_{d-1}R^{d+1}}
 \int_{\partial B_R}y\,u(y)\,\mathrm d\sigma(y).
\]
Wykorzystując $|y|=R$ oraz
$\sigma(\partial B_R)=\omega_{d-1}R^{d-1}$,
z nierówności trójkąta dla całki wektorowej otrzymujemy
\[
 |\nabla u(0)|\le
 \frac{d}{\omega_{d-1}R^{d+1}}\,
 R\,\omega_{d-1}R^{d-1}\sup_{\partial B_R}|u|
 =\frac dR\sup_{\partial B_R}|u|.
\]
Z ciągłości na domknięciu $\sup_{\partial B_R}|u|
\le\sup_{B_R}|u|$. Dla $x\in B_r(x_0)$ kula
$\overline{B_{(R-r)/2}(x)}$ jest zawarta w $B_R(x_0)$
z nierówności trójkąta. Stosując pierwsze oszacowanie
w tej kuli i biorąc supremum po $x$, dostajemy drugie
oszacowanie z tezy.
W istocie można wziąć nawet promień $R-r$, ponieważ
$|x-x_0|<r$; daje to stałą $d/(R-r)$ zamiast $2d/(R-r)$.

Na mocy twierdzenia~\ref{thm:pde7-weyl} funkcja harmoniczna
jest gładka, a każda jej pochodna jest harmoniczna, gdyż
pochodne komutują z laplasjanem. Dla wielowskaźnika
$\alpha$ o długości $m\ge1$ dzielimy odstęp $R-r$
na $m$ równych części i kładziemy
$r_j=R-j(R-r)/m$ dla $j=0,\ldots,m$.
Stosujemy drugie oszacowanie kolejno do pochodnych
rzędów $0,1,\ldots,m-1$ na kulach o promieniach $r_j$.
Każde przejście daje czynnik co najwyżej $dm/(R-r)$,
gdyż moduł pojedynczej składowej gradientu nie przekracza
jego normy euklidesowej. W rezultacie
\[
 \sup_{B_r(x_0)}|\partial^\alpha u|
 \le \left(\frac{dm}{R-r}\right)^m\sup_{B_R(x_0)}|u|.
\]
Jeśli supremum po prawej jest nieskończone, nierówność
rozumiemy w rozszerzonym sensie; na każdej kuli o domknięciu
zawartym w dziedzinie supremum jest skończone.
\end{proof}

\section{Subharmoniczność i nierówność wartości średniej}
\label{sec:pde14-subharmonic}
\begin{definicja}[Funkcja subharmoniczna w sensie klasycznym]
Niech $U\subset\R^d$ będzie otwarta, a $u\in C^2(U)$.
Nazywamy $u$ \emph{subharmoniczną}, jeśli $\Delta u\ge0$
w $U$, zaś \emph{superharmoniczną}, jeśli $\Delta u\le0$.
Istnieją również ogólniejsze pojęcia subharmoniczności
dla funkcji półciągłych z góry oraz superharmoniczności
dla funkcji półciągłych z dołu, związane z odpowiednimi
nierównościami dystrybucyjnymi. W tej sekcji twierdzenia
punktowe dowodzimy w klasie rzeczywistych funkcji $C^2$.
\end{definicja}

\begin{twierdzenie}[Nierówność wartości średniej]
\label{thm:pde14-sub-mean}
Jeżeli $u\in C^2(U)$ jest subharmoniczna i
$\overline{B_R(x_0)}\subset U$, to dla każdego $0<r\le R$
\[
 u(x_0)\le \frac{1}{\sigma(\partial B_r(x_0))}
 \int_{\partial B_r(x_0)}u(y)\,\mathrm d\sigma(y).
\]
Zachodzi również nierówność dla średniej po kuli:
\[
 u(x_0)\le \frac1{\lambda_d(B_r(x_0))}
 \int_{B_r(x_0)}u(y)\,\mathrm d\lambda_d(y).
\]
Dla funkcji superharmonicznej obie nierówności mają przeciwny zwrot.
\end{twierdzenie}
\begin{proof}
Niech $m(r)$ oznacza średnią sferyczną z tezy.
Korzystamy z rachunku przeprowadzonego w dowodzie
twierdzenia~\ref{thm:pde3-mean}, bez zakładania $\Delta u=0$.
Po podstawieniu $y=x_0+r\theta$, $\theta\in S^{d-1}$,
mamy $m(r)=\omega_{d-1}^{-1}
\int_{S^{d-1}}u(x_0+r\theta)\,\mathrm d\sigma(\theta)$.
Na zwartym otoczeniu przedziału promieni $[0,R]$
dopuszczalnym dzięki $\overline{B_R(x_0)}\subset U$
funkcja podcałkowa i jej pochodna względem $r$
są jednostajnie ograniczone na sferze. Twierdzenie
o wartości średniej i zbieżność zmajoryzowana
uzasadniają różniczkowanie pod całką. Otrzymujemy
\[
 m'(r)=\frac1{\omega_{d-1}}
 \int_{S^{d-1}}\nabla u(x_0+r\theta)\cdot\theta\,
 \mathrm d\sigma(\theta)
 =\frac1{\omega_{d-1}r^{d-1}}
 \int_{B_r(x_0)}\Delta u(y)\,\mathrm d\lambda_d(y).
\]
Druga równość wynika ze skalowania miary powierzchniowej
i twierdzenia Gaussa zastosowanego do pola $\nabla u$,
klasy $C^1$ w otoczeniu domkniętej kuli; normalna
zewnętrzna w punkcie $x_0+r\theta$ jest równa $\theta$.
Dla $\Delta u\ge0$ mamy $m'(r)\ge0$.
Z ciągłości $u$ w $x_0$ otrzymujemy
$|m(r)-u(x_0)|\le
\sup_{\theta\in S^{d-1}}|u(x_0+r\theta)-u(x_0)|\to0$,
gdy $r\downarrow0$. Monotoniczność $m$ daje pierwszą
nierówność. Całkując ją dla $0<s<r$ z wagą
$\omega_{d-1}s^{d-1}$, otrzymujemy
\[
 \int_{B_r(x_0)}u(y)\,\mathrm d\lambda_d(y)
 =\omega_{d-1}\int_0^r m(s)s^{d-1}\,\mathrm d\lambda_1(s)
 \ge u(x_0)\frac{\omega_{d-1}r^d}{d}.
\]
Ostatni współczynnik jest miarą kuli. Dzielimy przez niego.
Dla $\Delta u\le0$ stosujemy obie nierówności do $-u$.
\end{proof}

\begin{uwaga}[Wersja dystrybucyjna]
Dla $u\in L^1_{\mathrm{loc}}(U)$ warunek
$\Delta u\ge0$ w $\mathcal D'(U)$ oznacza
$\langle\Delta u,\varphi\rangle\ge0$ dla nieujemnych
$\varphi\in C_c^\infty(U)$. Dla nieujemnego standardowego
jądra wygładzającego $\rho_\varepsilon$ o całce $1$
i nośniku w $\overline{B_\varepsilon(0)}$ określamy
$u_\varepsilon=u*\rho_\varepsilon$ na
$U_\varepsilon=\{x\in U:\operatorname{dist}(x,\R^d\setminus U)>
\varepsilon\}$ (dla $U=\R^d$ przyjmujemy $U_\varepsilon=U$).
Z twierdzenia~\ref{thm:pde7-mollify} mamy
\[
 \Delta u_\varepsilon(x)
 =\langle\Delta u,\rho_\varepsilon(x-\cdot)\rangle\ge0,
 \qquad x\in U_\varepsilon,
\]
gdyż funkcja po prawej jest nieujemnym testem o zwartym
nośniku w $U$. Zatem $u_\varepsilon$ jest gładka
i subharmoniczna. Twierdzenie~\ref{thm:pde7-approx-id}
daje zbieżność do $u$ w $L^1_{\mathrm{loc}}(U)$.
Taka zbieżność nie wystarcza sama w sobie do przejścia
do nierówności punktowej dla dowolnego reprezentanta.
Pełna teoria wyróżnionych reprezentantów półciągłych
z góry (mogących przyjmować wartość $-\infty$)
wymaga dodatkowego twierdzenia, którego tutaj nie dowodzimy.
\end{uwaga}

\begin{twierdzenie}[Silna zasada minimum dla funkcji superharmonicznych]
\label{thm:pde14-strong-super}
Niech $U\subset\R^d$ będzie obszarem i niech
$u\in C^2(U)$ spełnia $\Delta u\le0$ w $U$.
Jeżeli $u$ osiąga minimum globalne względem $U$ w punkcie
$x_0\in U$, to jest stała na $U$.
\end{twierdzenie}
\begin{proof}
Niech $m=u(x_0)=\inf_Uu$. Z twierdzenia
~\ref{thm:pde14-sub-mean} w wersji superharmonicznej wynika
\[
 u(x_0)\ge\frac{1}{\sigma(\partial B_r(x_0))}
 \int_{\partial B_r(x_0)}u(y)\,\mathrm d\sigma(y)
\]
dla każdej kuli o domknięciu zawartym w $U$.
Ponieważ $u\ge m$, średnia wynosi jednocześnie co
najmniej $m$. Średnia funkcji nieujemnej $u-m$
na sferze jest więc zerowa. Gdyby w pewnym punkcie
sfery $u-m>0$, z ciągłości byłaby dodatnia, z dodatnią
dolną granicą, na fragmencie sfery o dodatniej mierze,
co przeczyłoby zerowej średniej. Stąd $u=m$ na każdej
takiej sferze, a więc w całej małej kuli o środku $x_0$.
Dla każdego $z$ ze zbioru $E=\{x\in U:u(x)=m\}$
powtarzamy argument z $z$ zamiast $x_0$. Zbiór $E$ jest
zatem otwarty w $U$, a dzięki ciągłości również względnie
domknięty. Jest niepusty, więc ze spójności $U$ mamy $E=U$.
Stosując wynik do $-u$, otrzymujemy analogiczną mocną
zasadę maksimum dla funkcji subharmonicznych.
\end{proof}

\begin{wniosek}[Lemat Hopfa dla nierówności $-\Delta u\ge0$]
\label{cor:pde14-hopf-super}
Niech $U\subset\R^d$ będzie obszarem ograniczonym,
$u\in C^2(U)\cap C(\overline U)$, $\Delta u\le0$ w $U$.
Załóżmy, że $u$ nie jest stała,
$u(x_0)=\min_{\overline U}u$ dla $x_0\in\partial U$
i istnieje kula $B_R(z)\subset U$ styczna do brzegu w $x_0$.
Wówczas, przy $\nu=(x_0-z)/R$,
\[
 \liminf_{t\downarrow0}\frac{u(x_0-t\nu)-u(x_0)}t>0.
\]
Jeśli istnieje skończona pochodna w kierunku normalnej
do kuli, obliczana od strony wnętrza zgodnie z twierdzeniem
~\ref{thm:pde14-hopf}, jest ona ujemna.
\end{wniosek}
\begin{proof}
Z twierdzenia~\ref{thm:pde14-strong-super} wynika,
że $u(x)>u(x_0)$ dla każdego $x\in U$. Odejmujemy
$u(x_0)$ i używamy tej samej bariery
$w(x)=e^{-\alpha|x-z|^2}-e^{-\alpha R^2}$,
co w dowodzie twierdzenia~\ref{thm:pde14-hopf}.
Na pierścieniu $R/2<|x-z|<R$ mamy $\Delta w>0$
dla $\alpha>2d/R^2$, a więc
$\Delta(u-\varepsilon w)\le-\varepsilon\Delta w<0$
dla $\varepsilon>0$. Po odjęciu minimum funkcja $u$
jest dodatnia na zwartej wewnętrznej sferze; oznaczamy
jej minimum przez $m>0$ i wybieramy
$0<\varepsilon\le m/(e^{-\alpha R^2/4}-e^{-\alpha R^2})$.
Wtedy $u-\varepsilon w\ge0$ na wewnętrznej sferze,
a na zewnętrznej również, gdyż $w=0$ i $u\ge0$.
Ujemne minimum na domkniętym pierścieniu musiałoby
być osiągnięte we wnętrzu, gdzie laplasjan byłby
nieujemny, sprzecznie z powyższą ścisłą nierównością.
Stąd $u\ge\varepsilon w$.
Dla $0<t<R/2$ podstawiamy $x=x_0-t\nu$ i otrzymujemy
\[
 \liminf_{t\downarrow0}
 \frac{u(x_0-t\nu)-u(x_0)}{t}
 \ge 2\varepsilon\alpha R e^{-\alpha R^2}>0,
\]
gdzie po lewej przywrócono odjętą wartość minimum.
Istnienie pochodnej z tezy daje jej ścisłą ujemność.
\end{proof}



\section{Nierówność Caccioppolego dla współczynników mierzalnych}
\label{sec:pde14-caccioppoli}
Przechodzimy ponownie do form dywergencyjnych, ale
bez założenia gładkości współczynników.

\begin{twierdzenie}[Oszacowanie wewnętrzne Caccioppolego]
\label{thm:pde14-caccioppoli}
Niech $\Omega\subset\R^d$ będzie otwarta,
$A\in L^\infty(\Omega;\R^{d\times d})$,
$\xi^{\mathsf T}A(x)\xi\ge\lambda|\xi|^2$ i
$|A(x)\xi|\le\Lambda|\xi|$ dla wszystkich $\xi\in\R^d$
na wspólnym zbiorze pełnej miary, gdzie $\lambda,\Lambda>0$.
Jeżeli rzeczywista funkcja $u\in H^1_{\mathrm{loc}}(\Omega)$
spełnia $-\diver(A\nabla u)=0$ dystrybucyjnie, to dla każdej
rzeczywistej $\eta\in C_c^\infty(\Omega)$ oraz $k\in\R$ zachodzi
\begin{equation}\label{eq:pde14-caccioppoli}
 \int_\Omega\eta^2|\nabla u|^2\,\mathrm d\lambda_d(x)
 \le \frac{4\Lambda^2}{\lambda^2}
 \int_\Omega (u-k)^2|\nabla\eta|^2\,\mathrm d\lambda_d(x).
\end{equation}
W szczególności, gdy $B_{2r}(x_0)\Subset\Omega$,
\[
 \int_{B_r(x_0)}|\nabla u|^2\,\mathrm d\lambda_d(x)
 \le \frac{C\Lambda^2}{\lambda^2r^2}
 \int_{B_{2r}(x_0)}|u-k|^2\,\mathrm d\lambda_d(x),
\]
gdzie $C$ zależy wyłącznie od wyboru standardowego obcięcia.
\end{twierdzenie}
\begin{proof}
Ponieważ operator nie ma członu zerowego rzędu,
$v=u-k$ spełnia to samo równanie i $\nabla v=\nabla u$.
Wybieramy otwarte $W\Subset\Omega$ zawierające
$\operatorname{supp}\eta$. Z lokalnej przynależności
$v\in H^1(W)$ i reguły iloczynu
~\ref{thm:pde9-products} mamy
$\varphi=\eta^2v\in H^1(W)$ oraz
$\nabla\varphi=\eta^2\nabla v+2\eta v\nabla\eta$.
Funkcja $\varphi$ znika poza zwartym podzbiorem $W$.
Jej przedłużenie przez zero na $\R^d$ należy do $H^1(\R^d)$:
dla testu na $\R^d$ mnożymy go przez gładkie obcięcie
równe $1$ w otoczeniu $\operatorname{supp}\eta$
i o nośniku w $W$, a następnie stosujemy tożsamość
słabej pochodnej na $W$. Otrzymujemy pochodne
przedłużone przez zero. Wygładzenie tego przedłużenia
przez splot daje testy o nośnikach w $W$ dla dostatecznie
małego parametru, zbieżne do $\varphi$ w $H^1$
(z twierdzeń~\ref{thm:pde7-mollify}
i~\ref{thm:pde7-approx-id} zastosowanych do funkcji
i jej pierwszych pochodnych na ustalonym zwartym
otoczeniu ich nośników).
Wobec tego $\varphi\in H_0^1(W)$.
Ponieważ $A\nabla u\in L^2(W)$, całka
$\int_W A\nabla u\cdot\nabla\psi\,\mathrm d\lambda_d(x)$
jest ciągła w normie $H^1(W)$. Przechodząc do granicy
z testami gładkimi, wolno więc użyć $\varphi$ w równaniu
dystrybucyjnym. Otrzymujemy
\[
 \int_\Omega\eta^2 A\nabla u\cdot\nabla u\,\mathrm d\lambda_d(x)
 =-2\int_\Omega\eta (u-k) A\nabla u\cdot\nabla\eta\,\mathrm d\lambda_d(x).
\]
Z eliptyczności, oszacowania normy operatorowej $A$
i nierówności Cauchy'ego--Schwarza otrzymujemy
$\lambda X^2\le2\Lambda XY$, gdzie
\[
 X^2=\int_\Omega\eta^2|\nabla u|^2\,\mathrm d\lambda_d(x),
 \qquad
 Y^2=\int_\Omega (u-k)^2|\nabla\eta|^2\,\mathrm d\lambda_d(x).
\]
Gdy $X=0$, teza jest oczywista. W przeciwnym razie
$X\le 2\Lambda Y/\lambda$, co po podniesieniu do kwadratu
daje~\eqref{eq:pde14-caccioppoli}.
Niech $\chi\in C^\infty(\R)$, $0\le\chi\le1$,
będzie równa $1$ dla $t\le1$ i równa $0$ dla
$t\ge3/2$. Kładziemy $\eta(x)=\chi(|x-x_0|/r)$.
Funkcja jest gładka również w środku (jest tam lokalnie
stała), ma nośnik w $\overline{B_{3r/2}(x_0)}$,
jest równa $1$ na $B_r(x_0)$ oraz
$|\nabla\eta|\le C_0/r$, gdzie $C_0=\|\chi'\|_\infty$.
Podstawienie do~\eqref{eq:pde14-caccioppoli} daje
drugie oszacowanie z $C=4C_0^2$.
W szczególności możemy wybrać $k=0$ lub średnią
$u_{B_{2r}(x_0)}$, uzyskując oszacowanie przez oscylację
w normie $L^2$ bez dodatkowego założenia na znak $u$.
\end{proof}

\begin{uwaga}[Dlaczego nie dowodzimy tu ogólnej nierówności Harnacka?]
Oszacowanie Caccioppolego jest punktem wyjścia do iteracji
De Giorgiego i Mosera. Pełna nierówność Harnacka dla
nieujemnych słabych rozwiązań przy mierzalnych współczynnikach
jest wynikiem znacznie głębszym niż nierówność dla funkcji
harmonicznych wyprowadzona z jądra Poissona. Przedstawimy
jej dowód w twierdzeniu~\ref{thm:pde17-harnack}
w rozdziale o regularności De Giorgiego--Nasha--Mosera;
nie używamy jej tutaj jako przesłanki.
\end{uwaga}

\section{Funkcja Greena i porównanie dla problemu Dirichleta}
\label{sec:pde14-green}
Niech $d\ge2$ oraz $\Phi$ będzie rozwiązaniem fundamentalnym
operatora $-\Delta$ w $\R^d$:
\[
 \Phi(z)=\begin{cases}
 \dfrac{|z|^{2-d}}{(d-2)\omega_{d-1}},&d\ge3,\\[4pt]
 -\dfrac1{2\pi}\log|z|,&d=2.
 \end{cases}
\]
W obu przypadkach $-\Delta\Phi=\delta_0$ w sensie
dystrybucyjnym. Nie utożsamiamy jądra Greena z samym
rozwiązaniem fundamentalnym: warunek brzegowy wymaga
odjęcia części harmonicznej.

\begin{definicja}[Funkcja Greena Dirichleta]
Niech $\Omega\subset\R^d$ będzie obszarem ograniczonym,
w którym dla każdego $y\in\Omega$ istnieje funkcja
$h_y\in C(\overline\Omega)\cap C^2(\Omega)$ harmoniczna
w $\Omega$ i taka, że $h_y(x)=\Phi(x-y)$ dla
$x\in\partial\Omega$. Definiujemy
$G_\Omega(x,y)=\Phi(x-y)-h_y(x)$ dla $x\ne y$.
Wówczas $-\Delta_xG_\Omega(x,y)=\delta_y$ dystrybucyjnie,
a $G_\Omega(\cdot,y)=0$ na brzegu.
\end{definicja}

\begin{twierdzenie}[Dodatniość funkcji Greena]
\label{thm:pde14-green-pos}
Przy powyższych założeniach dla każdego $y\in\Omega$
$G_\Omega(x,y)>0$ dla $x\in\Omega\setminus\{y\}$.
\end{twierdzenie}
\begin{proof}
Dla małego $\rho>0$ kula $\overline{B_\rho(y)}$ mieści
się w $\Omega$. Ponieważ $h_y$ jest ograniczona na
$\overline\Omega$, podczas gdy $\Phi(x-y)\to+\infty$
gdy $x\to y$ (również przy $d=2$), istnieje $\rho_0>0$
takie, że $G_\Omega(x,y)>0$ dla $|x-y|=\rho\le\rho_0$.
Na zewnętrznej granicy $\partial\Omega$ funkcja $G$ znika.
W obszarze pierścieniowym
$\Omega\setminus\overline{B_\rho(y)}$ jest harmoniczna;
słaba zasada minimum daje $G\ge0$ wszędzie w tym obszarze. Mocna zasada minimum daje ścisłą dodatniość.
Istotnie, $\Omega\setminus\overline{B_\rho(y)}$ jest spójna:
dowolną łamaną w $\Omega$ przecinającą kulę można zmodyfikować
wzdłuż sfery o promieniu nieco większym niż $\rho$,
która nadal leży w $\Omega$. Dowolny $x\ne y$
możemy objąć takim pierścieniem, wybierając $\rho<|x-y|$.
\end{proof}

\begin{lemat}[Jawna funkcja Greena kuli i jej pochodna normalna]
\label{lem:pde14-green-ball}
Niech $d\ge2$, $R>0$ i $x\in B_R\setminus\{0\}$.
Zapiszmy $x^*=R^2x/|x|^2$, a więc $|x^*|>R$. Dla
$y\in B_R$, $y\ne x$, funkcja Greena ma postać
\[
 G_{B_R}(x,y)=\begin{cases}
 \Phi(x-y)-(R/|x|)^{d-2}\Phi(x^*-y),&d\ge3,\\[3pt]
 \dfrac1{2\pi}\log\left(\dfrac{|x|\,|x^*-y|}
 {R|x-y|}\right),&d=2.
 \end{cases}
\]
W punkcie $x=0$ rozumiemy ją przez granicę. Jądro jest
symetryczne: $G_{B_R}(x,y)=G_{B_R}(y,x)$. Ponadto
$G_{B_R}(x,y)=0$ dla $|y|=R$, a na brzegu zachodzi
$-\partial_{\nu_y}G_{B_R}(x,y)=P_R(x,y)$, gdzie
$\nu_y=y/R$.
\end{lemat}
\begin{proof}
Dla $|y|=R$ otrzymujemy tożsamość geometryczną
$|x^*-y|=(R/|x|)|x-y|$. Po jej podstawieniu oba wzory
na $G$ przyjmują wartość zero. W zmiennej $y$ pierwszy
składnik jest harmoniczny poza biegunem $x$, drugi zaś
jest harmoniczny w całej kuli, ponieważ jego biegun
$x^*$ leży na zewnątrz. W pobliżu $y=x$ pierwszy
składnik jest rozwiązaniem fundamentalnym $-\Delta$,
więc $-\Delta_yG=\delta_x$.
Symetria wynika z równości
\[
 |x|^2|x^*-y|^2=R^4-2R^2x\cdot y+|x|^2|y|^2,
\]
której prawa strona jest symetryczna w $x,y$.

Dla $d\ge3$ wykorzystujemy
$\nabla_y\Phi(x-y)=(x-y)/(\omega_{d-1}|x-y|^d)$.
Różniczkując na zewnętrznej sferze i używając ponownie
$|x^*-y|=(R/|x|)|x-y|$, otrzymujemy
\begin{align*}
 \partial_{\nu_y}G(x,y)
 &=\frac{x\cdot y/R-R}{\omega_{d-1}|x-y|^d}
 -\frac{(|x|/R)^2(x^*\cdot y/R-R)}
 {\omega_{d-1}|x-y|^d}\\
 &=\frac{|x|^2-R^2}{\omega_{d-1}R|x-y|^d}
 =-P_R(x,y).
\end{align*}
Dla $d=2$ pochodna logarytmu spełnia tę samą
formułę, gdyż $\omega_1=2\pi$ i wykładnik $d-2$ wynosi
zero; składnik stały względem $y$ nie wnosi pochodnej.
Przejście $x\to0$ daje pozostały przypadek.
\end{proof}

\begin{twierdzenie}[Reprezentacja Greena w kuli]
\label{thm:pde14-green-repr}
Niech $d\ge2$, $R>0$, $f\in C(\overline{B_R})$,
$g\in C^2(\partial B_R)$, a
$u\in C^2(B_R)\cap C^1(\overline{B_R})$ spełnia
$-\Delta u=f$ w $B_R$ i $u=g$ na $\partial B_R$.
Wówczas dla $x\in B_R$ zachodzi
\begin{equation}\label{eq:pde14-green-repr}
 u(x)=\int_{B_R}G_{B_R}(x,y)f(y)\,\mathrm d\lambda_d(y)
 +\int_{\partial B_R}P_R(x,y)g(y)\,\mathrm d\sigma(y).
\end{equation}
W szczególności dla $f\ge0$ i $g\ge0$ mamy $u\ge0$.
\end{twierdzenie}
\begin{proof}
Dla ustalonego $x$ stosujemy drugą tożsamość Greena
w $B_R\setminus\overline{B_\rho(x)}$ do pary
$u(y)$ i $G_{B_R}(x,y)$. W tym obszarze $G$ jest
harmoniczna i $\Delta u=-f$. Otrzymujemy
\[
 \int_{B_R\setminus B_\rho(x)}G_{B_R}(x,y)f(y)
 \,\mathrm d\lambda_d(y)
 =\int_{\partial B_R}g\,\partial_{\nu_y}G\,
 \mathrm d\sigma+I_\rho,
\]
gdzie $I_\rho$ to wyraz z małej sfery, z normalną
skierowaną do jej środka. Ponieważ
$G(x,y)=\Phi(x-y)+O(1)$ w pobliżu $y=x$ oraz
$\partial_{\nu}G=\partial_{\nu}\Phi+O(1)$,
bezpośrednie obliczenie strumienia $\Phi$ daje
$\lim_{\rho\downarrow0}I_\rho=u(x)$.
Ponadto całka objętościowa przechodzi do granicy:
osobliwość $\Phi$ jest lokalnie całkowalna dla $d\ge2$.
Tożsamość $-\partial_{\nu_y}G_{B_R}(x,y)=P_R(x,y)$
została jawnie wyprowadzona w lemacie~\ref{lem:pde14-green-ball}.
Po przeniesieniu wyrazów otrzymujemy~\eqref{eq:pde14-green-repr}.
Wniosek o znaku wynika z dodatniości obu jąder.
\end{proof}

\begin{uwaga}[Uwaga o znakach we wzorze Greena]
W powyższym dowodzie normalna na \emph{wewnętrznej}
sferze $\partial B_\rho(x)$ jest skierowana \emph{do środka}
wyciętej kuli. To właśnie dlatego wyraz z tej sfery
ma znak prowadzący do $+u(x)$ w reprezentacji.
Rozpisanie orientacji na każdej granicy chroni przed
częstą pomyłką znaku przy przejściu od $\Phi$ do $G$.
\end{uwaga}

\section{Rysunek: bariera przy brzegu i geometria nierówności}
\label{sec:pde14-figure}
\begin{center}
\begin{tikzpicture}[scale=0.88,>=Latex]
  % Dwie kule styczne wewnetrznie: x0 nalezy do obu brzegow.
  \fill[blue!5] (-0.6,0) circle (2.45);
  \draw[thick,blue!65!black] (-0.6,0) circle (2.45);
  \draw[thick,orange!90!black] (0,0) circle (1.85);
  \fill (0,0) circle(2pt) node[below left] {$z$};
  \fill (1.85,0) circle(2.5pt) node[above right] {$x_0$};
  \draw[dashed] (0,0)--(1.85,0) node[midway,above] {$R$};
  \draw[-{Latex[length=2.5mm]},very thick] (1.85,0) -- (2.65,0)
    node[right] {$\nu$};
  \draw[-{Latex[length=2.5mm]},very thick,red!75!black]
    (1.78,-0.28) -- (0.90,-0.28);
  \node[red!70!black,below] at (0.15,-0.65) {kierunek $-\nu$};
  \node[blue!70!black] at (-2.3,1.25) {$\Omega$};
  \node[orange!85!black] at (-0.7,1.20) {$B_R(z)$};
\end{tikzpicture}

{\small Geometria lematu Hopfa: kula wewnętrzna dotyka brzegu
$\partial\Omega$ w $x_0$, a bariera wykrywa dodatni przyrost
w kierunku do wnętrza $-\nu$.}
\end{center}

\section{Podsumowanie i przejście do teorii regularności}
\label{sec:pde14-summary}
Wykazaliśmy słabą zasadę maksimum przy mierzalnych
współczynnikach i $c\ge0$, nie wykorzystując gładkości
ani symetrii części głównej. Dla laplasjanu udowodniliśmy
mocną zasadę maksimum, lemat Hopfa, ilościową nierówność
Harnacka, oszacowania gradientu i własności subharmoniczne.
Nierówność Caccioppolego pozostaje prawdziwa dla słabych
rozwiązań z mierzalną macierzą współczynników, ale sama
nie zapewnia od razu ciągłości rozwiązań ani nierówności
Harnacka. Te głębsze zagadnienia prowadzą do
teorii De Giorgiego--Nasha--Mosera oraz regularności
eliptycznej w kolejnych rozdziałach.

% Iteracja 15: regularność eliptyczna, całość rozdziału.
